\documentclass[11pt]{article}

\usepackage[T1]{fontenc}
\usepackage{lmodern}
\usepackage{amsmath,amssymb,amsthm}
\usepackage{booktabs}
\usepackage[margin=1in]{geometry}
\usepackage{microtype}
\usepackage[numbers,sort&compress]{natbib}
\usepackage[hidelinks]{hyperref}
\hypersetup{
  pdftitle={Exact Cut-Uniform Ghost-Penalty Thresholds in the Aligned Two-Cell CutFEM Model for Degrees One Through Five},
  pdfauthor={Xiaojian Zeng}
}

\newtheorem{theorem}{Theorem}
\newtheorem{lemma}{Lemma}
\newcommand{\R}{\mathbb R}
\newcommand{\PP}{\mathbb P}
\newcommand{\eps}{\varepsilon}

\title{Exact Cut-Uniform Ghost-Penalty Thresholds in the Aligned
Two-Cell CutFEM Model for Degrees One Through Five}
\author{Xiaojian Zeng\\
\small PBC School of Finance, Tsinghua University, 100190\\
\small \texttt{xiaojianzeng@foxmail.com}}
\date{}

\begin{document}
\maketitle

\begin{abstract}
For the aligned two-element CutFEM model of Wichrowski, with polynomial
degree $p\in\{1,2,3,4,5\}$ and Nitsche parameter
$\gamma_D=5p(p+1)$, we determine the least cut-uniform ghost weight.
It is
\[
 \Gamma_p=\frac{1}{p(4p+5)};
\]
hence the five values are $1/9,1/26,1/51,1/84,1/125$.  The supremum is
not attained at an interior cut fraction: in every degree its unique
limiting maximizer is $\kappa\to0^+$, while $\kappa\to1^-$ is not
maximizing.  The proof is an exact congruence reduction.  At the claimed
weight the remaining Schur matrix has fraction-free leading minors whose
polynomial coefficients have certified positive signs.  A rank-one
downdate at the vanishing cut proves necessity.  This resolves the finite
two-cell question; it does not determine the stability threshold of an
assembled CutFEM operator.
\end{abstract}

\section{Context and result}

Ghost penalties were introduced to make unfitted finite element methods
robust with respect to the position of the physical boundary relative to
the background mesh \citep{Burman2010}; for CutFEM context see
\citet{BurmanEtAl2015}.  More recent designs relate stabilized functionals
to discrete finite element extensions \citep{BurmanHansboLarson2026}.
Wichrowski's aligned two-cell model isolates an exact finite-dimensional
question inside this setting \citep[Proposition~5 and Appendix~B]
{Wichrowski2026}.  That source proved the vanishing-cut threshold and a
global lower bound, but left equality with the full cut-uniform threshold
unproved.  Exact rational sign certificates have also been used for other
parameterized Hermitian positivity problems \cite{ZengEtAl2026}; the
certificate here is specialized to the polynomial matrices of the
two-cell form.

Let $\PP_p$ denote the real polynomials of degree at most $p$.  For
$0<\kappa<1$, $\gamma>0$, and $(u,q)\in\PP_p^2$, set
\begin{align}
a^{(2)}_{\kappa,\gamma}(u,q)
={}&\int_{-1}^{0}(u')^2\,dx+\int_0^\kappa(q')^2\,dx
-2q'(\kappa)q(\kappa)+Dq(\kappa)^2 \notag\\
&+\gamma\sum_{r=1}^p
 \frac{\bigl(q^{(r)}(0)-u^{(r)}(0)\bigr)^2}{(r!)^2},
 \qquad D:=5p(p+1).                                      \label{eq:form}
\end{align}
This is exactly the dimensionless one-dimensional reduction in
\citet[Eq.~(60)]{Wichrowski2026}.  Define
\[
 \tau_p(\kappa):=\inf\{\gamma>0:
 a^{(2)}_{\kappa,\gamma}\text{ is positive semidefinite on }\PP_p^2\},
 \qquad
 \Gamma_p:=\sup_{0<\kappa<1}\tau_p(\kappa).
\]

\begin{theorem}\label{thm:main}
For every $p\in\{1,2,3,4,5\}$,
\[
 \Gamma_p=\frac1{D-p^2}=\frac1{p(4p+5)}.
\]
For every fixed $\kappa\in(0,1)$ one has
$\tau_p(\kappa)<\Gamma_p$.  Moreover,
\[
 \lim_{\kappa\to0^+}\tau_p(\kappa)=\Gamma_p,
 \qquad
 \limsup_{\kappa\to1^-}\tau_p(\kappa)<\Gamma_p.
\]
Thus there is no interior maximizer and the only one-sided limiting
maximizer is $\kappa\to0^+$.
\end{theorem}

\section{Congruence reduction}

Write
\[
 u(x)=u_0+\sum_{i=1}^p\alpha_i x^i,
 \qquad q(x)=q_0+\sum_{i=1}^p\beta_i x^i,
\]
and put $\alpha=(\alpha_1,\ldots,\alpha_p)^T$ and
$\beta=(\beta_1,\ldots,\beta_p)^T$.  The coefficient $u_0$ is an
unused null direction.  Replacing $q_0$ by the independent coordinate
$y=q(\kappa)$ is invertible for every $\kappa$.  Define
\begin{equation}
 U_{ij}=\frac{ij(-1)^{i+j-2}}{i+j-1},\qquad
 K_{ij}(\kappa)=\frac{ij\kappa^{i+j-1}}{i+j-1},\qquad
 v_i(\kappa)=i\kappa^{i-1}.                              \label{eq:UKv}
\end{equation}
Direct integration, using the factorial normalization in
\eqref{eq:form}, and completion of the boundary square give the exact
identity
\begin{align}
a^{(2)}_{\kappa,\gamma}
={}&D\left(y-\frac{v^T\beta}{D}\right)^2
+\alpha^TU\alpha+\gamma\|\beta-\alpha\|^2
+\beta^T\left(K-\frac{vv^T}{D}\right)\beta .             \label{eq:firstsquare}
\end{align}
The matrix $U$ is the Gram matrix of the derivative monomials
$ix^{i-1}$ on $(-1,0)$ and is therefore positive definite.

Set
\[
 g_*:=\frac1{D-p^2},\qquad P:=U+g_*I,
\]
so $P\succ0$.  A second exact completion gives
\begin{align}
\alpha^TU\alpha+g_*\|\beta-\alpha\|^2
+\beta^T\left(K-\frac{vv^T}{D}\right)\beta
={}&(\alpha-g_*P^{-1}\beta)^TP(\alpha-g_*P^{-1}\beta)
+\beta^TS_p(\kappa)\beta,                                \label{eq:secondsquare}\\
S_p(\kappa):={}&K(\kappa)-\frac{v(\kappa)v(\kappa)^T}{D}
+g_*I-g_*^2P^{-1}.                                        \label{eq:schur}
\end{align}

\begin{lemma}[Exact Schur certificates]\label{lem:cert}
For $p=1,\ldots,5$, $S_p(\kappa)\succeq0$ for every
$\kappa\ge0$ and $S_p(\kappa)\succ0$ for every $\kappa>0$.
\end{lemma}

\begin{proof}
For $p=1$, direct substitution in \eqref{eq:schur} gives
$S_1(\kappa)=\kappa$.  For $p=2$, $g_*=1/26$ and exact arithmetic gives
\[
S_2(\kappa)=
\begin{pmatrix}
\dfrac{2870\kappa+1}{2870}&
\dfrac{4305\kappa^2-287\kappa-15}{4305}\\[5pt]
\dfrac{4305\kappa^2-287\kappa-15}{4305}&
\dfrac{2(2870\kappa^3-287\kappa^2+75)}{4305}
\end{pmatrix}
\]
and
\[
 \det S_2(\kappa)=
 \frac{\kappa(4305\kappa^3+6\kappa^2+32\kappa+444)}{12915}.
\]
Its first leading minor is positive for $\kappa\ge0$, while its
determinant is nonnegative there and positive for $\kappa>0$.

For $p=3,4,5$, let $\Delta_{p,j}$ be the leading $j\times j$ principal
minor of $S_p$.  The exact coefficient data in Appendix~\ref{app:data}
have the form
\[
 \Delta_{p,j}(\kappa)=d_{p,j}^{-1}
 \sum_{r=0}^{j^2}c_{p,j,r}\kappa^r,
 \qquad d_{p,j}>0.
\]
For $j<p$, every listed coefficient is strictly positive.  For $j=p$,
the constant coefficient is zero and every coefficient of positive
degree is strictly positive.  Fraction-free Gaussian elimination
therefore has positive first $p-1$ pivots and a nonnegative last pivot
on $\kappa\ge0$; the last is positive exactly when $\kappa>0$.
This proves the two assertions by congruence.
\end{proof}

Equations \eqref{eq:firstsquare}--\eqref{eq:secondsquare} and
Lemma~\ref{lem:cert} show that the form at $g_*$ is positive
semidefinite for the full interval $0\le\kappa\le1$.  On all effective
coordinates it is positive definite when $\kappa>0$.

\section{Endpoint necessity and maximizers}

It remains to prove that the preceding upper bound is sharp.  At
$\kappa=0$ and an arbitrary $g>0$, complete the $\alpha$-square in
\eqref{eq:firstsquare}.  The remaining $\beta$-matrix is
\begin{equation}
 S_{p,g}(0)=A_g-\frac{e_1e_1^T}{D},\qquad
 A_g=gI-g^2(U+gI)^{-1}=gU(U+gI)^{-1}\succ0,               \label{eq:downdate}
\end{equation}
where $e_1=(1,0,\ldots,0)^T$.  Since
\begin{equation}
 A_g^{-1}=g^{-1}I+U^{-1},                                 \label{eq:ainv}
\end{equation}
the rank-one downdate criterion says that \eqref{eq:downdate} is
positive semidefinite exactly when
\begin{equation}
 g^{-1}+(U^{-1})_{11}\le D.                              \label{eq:rankonecriterion}
\end{equation}

For completeness, $(U^{-1})_{11}=p^2$ follows without a matrix
inversion.  It is the squared norm of endpoint evaluation on
$\PP_{p-1}$ in $L^2(-1,0)$.  The shifted Legendre basis has endpoint
values with squared norms $2m+1$, and hence its reproducing kernel at
the endpoint is
\[
 \sum_{m=0}^{p-1}(2m+1)=p^2.
\]
Thus \eqref{eq:rankonecriterion} is equivalent to $g\ge g_*$.  If
$g<g_*$, choose $\beta=A_g^{-1}e_1$ and take $\alpha,y$ to minimize the
two completed squares.  With $t=e_1^TA_g^{-1}e_1>D$, the remaining
value is the exact negative number
\[
 t-\frac{t^2}{D}<0.
\]
Continuity in $\kappa$ preserves this negative direction for all
sufficiently small positive cuts.  Since this holds for every $g<g_*$,
while Lemma~\ref{lem:cert} gives the global upper bound,
\[
 \lim_{\kappa\to0^+}\tau_p(\kappa)=g_*,
 \qquad \Gamma_p=g_*.
\]

For every fixed $\kappa>0$, the effective matrix at $g_*$ is positive
definite.  Its entries depend continuously on $(\kappa,g)$, so $g$ can
be decreased slightly while retaining positive semidefiniteness.  Hence
$\tau_p(\kappa)<g_*$ at every interior cut.  The same strictness at
$\kappa=1$, together with compact-neighborhood continuity, gives a
single weight below $g_*$ valid throughout a one-sided neighborhood of
$1$.  Therefore
$\limsup_{\kappa\to1^-}\tau_p(\kappa)<g_*$, completing the proof of
Theorem~\ref{thm:main}.

\section*{Computational and AI assistance disclosure}

Exact symbolic arithmetic in Python~3/SymPy was used to reconstruct
\eqref{eq:firstsquare}--\eqref{eq:schur}, to compute the fraction-free
leading minors, and to check every integer in Appendix~\ref{app:data}.
All computations used rational arithmetic; floating-point evidence is
not part of the proof.  OpenAI Codex assisted with symbolic checking,
proof auditing, and manuscript preparation.  A separate model-assisted
audit recomputed the decisive identities, but this is not a formal proof
assistant check or external peer review.  The author reviewed the
arguments and retains responsibility for the article.

\appendix
\section{Leading-minor coefficient certificate}\label{app:data}

For each line below, the positive integer before the semicolon is
$d_{p,j}$ and the parenthesized vector is
$(c_{p,j,0},c_{p,j,1},\ldots)$; omitted trailing coefficients are zero.
All entries are integers, so the sign verification is exact.

\begingroup\scriptsize
\allowdisplaybreaks
\paragraph{Degree $3$.}
\begin{align*}
(d_{3,1};\mathbf c_{3,1})={}&(288030;\ 307,288030),\\
(d_{3,2};\mathbf c_{3,2})={}&(5184540;\ 25,61422,23211,7368,1728180),\\
(d_{3,3};\mathbf c_{3,3})={}&(34563600;\ 0,5880,2840,2120,181635,143000,\\
&\hspace{35mm}310954,46422,5526,576060).
\end{align*}

\paragraph{Degree $4$.}
\begin{align*}
(d_{4,1};\mathbf c_{4,1})={}&(48007300;\ 51771,48007300),\\
(d_{4,2};\mathbf c_{4,2})={}&(3600547500;\ 24375,27996320,7713468,5177100,1200182500),\\
(d_{4,3};\mathbf c_{4,3})={}&(72010950000;\ 875,1739100,975540,1337000,114382590,\\
&\hspace{31mm}258434400,436320220,46280808,11648475,1200182500),\\
(d_{4,4};\mathbf c_{4,4})={}&(630095812500;\ 0,61000,62400,94500,6815025,23373000,\\
&\hspace{25mm}48129900,42105980,24482475,103345560,146389600,\\
&\hspace{25mm}175877940,193087755,135333240,7713468,1035420,60009125).
\end{align*}

\paragraph{Degree $5$.}
\begin{align*}
(d_{5,1};\mathbf c_{5,1})={}&(6916071394725;\ 6278291156,6916071394725),\\
(d_{5,2};\mathbf c_{5,2})={}&(20748214184175;\ 84567189,106874255547,28461576147,\\
&\hspace{39mm}25113164624,6916071394725),\\
(d_{5,3};\mathbf c_{5,3})={}&(414964283683500;\ 2668344,4810393332,2672153712,4198766224,\\
&\hspace{24mm}474811709205,1020581506614,1660048955862,\\
&\hspace{24mm}170769456882,56504620404,6916071394725),\\
(d_{5,4};\mathbf c_{5,4})={}&(72618749644612500;\ 333396,832077540,684736080,1293181680,\\
&\hspace{20mm}122210625915,462292088538,1214092038530,1612891692510,\\
&\hspace{20mm}1347981905385,3856221165605,10764841802388,\\
&\hspace{20mm}17921277398340,15187050676180,10287912491580,\\
&\hspace{20mm}569231522940,100452658496,6916071394725),\\
(d_{5,5};\mathbf c_{5,5})={}&(5123978974923858000;\ 0,32450544,55714176,139878144,12001478076,\\
&\hspace{15mm}85474806984,315059491656,646065363384,811064107476,\\
&\hspace{15mm}1531206650484,4815423196416,10864703913504,16644420401196,\\
&\hspace{15mm}19355277131064,17235330280896,12474909886544,\\
&\hspace{15mm}10263697225341,9027379754694,9240770097776,11347738287144,\\
&\hspace{15mm}8917299321651,4238179481074,1640073302706,56923152294,\\
&\hspace{15mm}6278291156,276642855789).
\end{align*}
\endgroup

\bibliographystyle{plainnat}
\bibliography{references}

\end{document}
